% GENERATED FILE. Edit canonical lessons, then run npm run generate:books.
% canonical-source-hash: fnv1a64:0b7eb090df910dca
% canonical-lessons: TE-C472-vanyaprani, TE-C472-abhayaranyam, TE-C472-vetadu, TE-C472-antarincipovu, TE-C472-valasapovu

\chapter{Wildlife}
\label{ch:te-wildlife}
\hlchaptermodality{472}

\begin{chapteropening}
\textbf{By the end of this chapter:} \emph{I can say a wild animal, a wildlife sanctuary, to hunt, to die out and to migrate.}

Complete the last lesson of chapter 472: I can say a wild animal, a wildlife sanctuary, to hunt, to die out and to migrate.
\end{chapteropening}

\section[\texorpdfstring{\te{వన్యప్రాణి} (vanyaprāṇi)}{వన్యప్రాణి (vanyaprāṇi)}]{\te{వన్యప్రాణి} (vanyaprāṇi) — a wild animal}

\label{lesson:TE-C472-vanyaprani}



\begin{quote}
Before the new one: say the Telugu for to damage, then the Telugu for to conserve.
\end{quote}

\subsection*{You'll want to know: \te{వన్యప్రాణి}}
\textbf{\te{వన్యప్రాణి}} — \emph{vanyaprāṇi} — \textquotedblleft{}a wild animal\textquotedblright{}.

\begin{quote}
\te{అడవిలో} \te{వన్యప్రాణులను} \te{చూడడానికి} \te{వెళ్ళాం}.
\end{quote}

\begin{quote}
\emph{aḍavilō vanyaprāṇulanu cūḍaḍāniki veḷḷāṁ.}
\end{quote}

\begin{quote}
We went to the forest to watch wild animals.
\end{quote}

Plural: \te{వన్యప్రాణులు}, \emph{vanyaprāṇulu}.

\te{వన్య}, of the forest, + \te{ప్రాణి}, a living being. Plural \te{వన్యప్రాణులు}, wildlife.

\begin{cousinweb}
\textbf{\te{వన్యప్రాణి}} (\emph{vanyaprāṇi}) is a compound of two words borrowed from Sanskrit: \emph{vanya}, \textquotedblleft{}growing in a forest\textquotedblright{} (CDIAL 11277), and \emph{prāṇin}, \textquotedblleft{}a breathing, living creature\textquotedblright{} (Monier-Williams).
\end{cousinweb}

\begin{grammarlens}[title={What you've built}]
One more word for getting around.
\end{grammarlens}

\subsection*{Your turn}
\begin{itemize}
\raggedright
  \item \emph{Say it:} \emph{vanyaprāṇi}
  \item \emph{Say it:} \emph{vanyaprāṇi}, once more
  \item \emph{Say it:} \emph{aḍavilō vanyaprāṇulanu cūḍaḍāniki veḷḷāṁ.} — We went to the forest to watch wild animals.
  \item \emph{Recall:} say the Telugu for to damage, then the Telugu for to conserve, then say \emph{vanyaprāṇi} again
\end{itemize}

\begin{tcolorbox}[breakable,colback=teal!4,colframe=teal!35!black,title={Before you move on}]
What does \te{వన్యప్రాణి} mean? (\textquotedblleft{}A wild animal\textquotedblright{}.) Say it once more.
\end{tcolorbox}
\section[\texorpdfstring{\te{అభయారణ్యం} (abhayāraṇyaṁ)}{అభయారణ్యం (abhayāraṇyaṁ)}]{\te{అభయారణ్యం} (abhayāraṇyaṁ) — a wildlife sanctuary}

\label{lesson:TE-C472-abhayaranyam}



\begin{quote}
Before the new one: say the Telugu for to conserve, then the Telugu for a wild animal.
\end{quote}

\subsection*{You'll want to know: \te{అభయారణ్యం}}
\textbf{\te{అభయారణ్యం}} — \emph{abhayāraṇyaṁ} — \textquotedblleft{}a wildlife sanctuary\textquotedblright{}.

\begin{quote}
\te{ఆ} \te{అభయారణ్యంలో} \te{పులులు} \te{ఉన్నాయి}.
\end{quote}

\begin{quote}
\emph{ā abhayāraṇyaṁlō pululu unnāyi.}
\end{quote}

\begin{quote}
There are tigers in that sanctuary.
\end{quote}

\te{అభయ}, freedom from fear, + \te{అరణ్యం}, a forest: a forest where animals need not fear.

\begin{cousinweb}
\textbf{\te{అభయారణ్యం}} (\emph{abhayāraṇyaṁ}) is a compound of two words borrowed from Sanskrit: \emph{abhaya}, \textquotedblleft{}fearlessness, safety\textquotedblright{}, and \emph{araṇya}, \textquotedblleft{}a wild place, a forest\textquotedblright{} (CDIAL 600).
\end{cousinweb}

\begin{grammarlens}[title={What you've built}]
One more word for getting around.
\end{grammarlens}

\subsection*{Your turn}
\begin{itemize}
\raggedright
  \item \emph{Say it:} \emph{abhayāraṇyaṁ}
  \item \emph{Say it:} \emph{abhayāraṇyaṁ}, once more
  \item \emph{Say it:} \emph{ā abhayāraṇyaṁlō pululu unnāyi.} — There are tigers in that sanctuary.
  \item \emph{Recall:} say the Telugu for to conserve, then the Telugu for a wild animal, then say \emph{abhayāraṇyaṁ} again
\end{itemize}

\begin{tcolorbox}[breakable,colback=teal!4,colframe=teal!35!black,title={Before you move on}]
What does \te{అభయారణ్యం} mean? (\textquotedblleft{}A wildlife sanctuary\textquotedblright{}.) Say it once more.
\end{tcolorbox}
\section[\texorpdfstring{\te{వేటాడు} (vēṭāḍu)}{వేటాడు (vēṭāḍu)}]{\te{వేటాడు} (vēṭāḍu) — to hunt}

\label{lesson:TE-C472-vetadu}



\begin{quote}
Before the new one: say the Telugu for a wild animal, then the Telugu for a wildlife sanctuary.
\end{quote}

\subsection*{You'll want to know: \te{వేటాడు}}
\textbf{\te{వేటాడు}} — \emph{vēṭāḍu} — \textquotedblleft{}to hunt\textquotedblright{}.

\begin{quote}
\te{ఇప్పుడు} \te{పులులను} \te{వేటాడకూడదు}.
\end{quote}

\begin{quote}
\emph{ippuḍu pululanu vēṭāḍakūḍadu.}
\end{quote}

\begin{quote}
Tigers must not be hunted now.
\end{quote}

Past (I …): \te{వేటాడాను}, \emph{vēṭāḍānu}.

\te{వేట}, a hunt, + \te{ఆడు}, to play, to do. The hunter is \te{వేటగాడు}.

\begin{cousinweb}
\textbf{\te{వేటాడు}} (\emph{vēṭāḍu}) is native Dravidian: DEDR 5527 lists Telugu \emph{vēṭa}, \textquotedblleft{}hunting, the chase\textquotedblright{}, and \emph{vēṭāḍu}, \textquotedblleft{}to hunt\textquotedblright{}.
\end{cousinweb}

\begin{grammarlens}[title={What you've built}]
One more word for getting around.
\end{grammarlens}

\subsection*{Your turn}
\begin{itemize}
\raggedright
  \item \emph{Say it:} \emph{vēṭāḍu}
  \item \emph{Say it:} \emph{vēṭāḍu}, once more
  \item \emph{Say it:} \emph{ippuḍu pululanu vēṭāḍakūḍadu.} — Tigers must not be hunted now.
  \item \emph{Recall:} say the Telugu for a wild animal, then the Telugu for a wildlife sanctuary, then say \emph{vēṭāḍu} again
\end{itemize}

\begin{tcolorbox}[breakable,colback=teal!4,colframe=teal!35!black,title={Before you move on}]
What does \te{వేటాడు} mean? (\textquotedblleft{}To hunt\textquotedblright{}.) Say it once more.
\end{tcolorbox}
\section[\texorpdfstring{\te{అంతరించిపోవు} (antariñcipōvu)}{అంతరించిపోవు (antariñcipōvu)}]{\te{అంతరించిపోవు} (antariñcipōvu) — to die out; to vanish}

\label{lesson:TE-C472-antarincipovu}



\begin{quote}
Before the new one: say the Telugu for a wildlife sanctuary, then the Telugu for to hunt.
\end{quote}

\subsection*{You'll want to know: \te{అంతరించిపోవు}}
\textbf{\te{అంతరించిపోవు}} — \emph{antariñcipōvu} — \textquotedblleft{}to die out; to vanish\textquotedblright{}.

\begin{quote}
\te{చాలా} \te{రకాల} \te{పక్షులు} \te{అంతరించిపోయాయి}.
\end{quote}

\begin{quote}
\emph{cālā rakāla pakṣulu antariñcipōyāyi.}
\end{quote}

\begin{quote}
Many kinds of birds have died out.
\end{quote}

Past (it …): \te{అంతరించిపోయింది}, \emph{antariñcipōyindi}.

Said of species, languages and old customs that disappear for good. Its subject is a thing, so the past shown is the third person: \te{అంతరించిపోయింది}, it died out.

\begin{cousinweb}
\textbf{\te{అంతరించిపోవు}} (\emph{antariñcipōvu}) is a Telugu verb built from Sanskrit \emph{antarita}, \textquotedblleft{}gone within, hidden, vanished\textquotedblright{} (Monier-Williams), with -\te{ఇంచు} and the native \te{పోవు}, \textquotedblleft{}to go\textquotedblright{} (DEDR 4572).
\end{cousinweb}

\begin{grammarlens}[title={What you've built}]
One more word for getting around.
\end{grammarlens}

\subsection*{Your turn}
\begin{itemize}
\raggedright
  \item \emph{Say it:} \emph{antariñcipōvu}
  \item \emph{Say it:} \emph{antariñcipōvu}, once more
  \item \emph{Say it:} \emph{cālā rakāla pakṣulu antariñcipōyāyi.} — Many kinds of birds have died out.
  \item \emph{Recall:} say the Telugu for a wildlife sanctuary, then the Telugu for to hunt, then say \emph{antariñcipōvu} again
\end{itemize}

\begin{tcolorbox}[breakable,colback=teal!4,colframe=teal!35!black,title={Before you move on}]
What does \te{అంతరించిపోవు} mean? (\textquotedblleft{}To die out; to vanish\textquotedblright{}.) Say it once more.
\end{tcolorbox}
\section[\texorpdfstring{\te{వలసపోవు} (valasapōvu)}{వలసపోవు (valasapōvu)}]{\te{వలసపోవు} (valasapōvu) — to migrate}

\label{lesson:TE-C472-valasapovu}



\begin{quote}
Before the new one: say the Telugu for to hunt, then the Telugu for to die out.
\end{quote}

\subsection*{You'll want to know: \te{వలసపోవు}}
\textbf{\te{వలసపోవు}} — \emph{valasapōvu} — \textquotedblleft{}to migrate\textquotedblright{}.

\begin{quote}
\te{చలికాలంలో} \te{చాలా} \te{పక్షులు} \te{దక్షిణానికి} \te{వలసపోతాయి}.
\end{quote}

\begin{quote}
\emph{calikālaṁlō cālā pakṣulu dakṣiṇāniki valasapōtāyi.}
\end{quote}

\begin{quote}
In winter many birds migrate south.
\end{quote}

Past (I …): \te{వలసపోయాను}, \emph{valasapōyānu}.

\te{వలస}, migration, + \te{పోవు}. Also said of villagers who move to the cities for work: \te{పనుల} \te{కోసం} \te{వలసపోయారు}.

\begin{cousinweb}
\textbf{\te{వలసపోవు}} (\emph{valasapōvu}) is native Dravidian: \te{వలస}, \textquotedblleft{}emigration, migration\textquotedblright{} (DEDR 5278), with \te{పోవు}, \textquotedblleft{}to go\textquotedblright{} (DEDR 4572).
\end{cousinweb}

\begin{grammarlens}[title={What you've built}]
One more word for getting around.
\end{grammarlens}

\subsection*{Your turn}
\begin{itemize}
\raggedright
  \item \emph{Say it:} \emph{valasapōvu}
  \item \emph{Say it:} \emph{valasapōvu}, once more
  \item \emph{Say it:} \emph{calikālaṁlō cālā pakṣulu dakṣiṇāniki valasapōtāyi.} — In winter many birds migrate south.
  \item \emph{Recall:} say the Telugu for to hunt, then the Telugu for to die out, then say \emph{valasapōvu} again
\end{itemize}

\begin{tcolorbox}[breakable,colback=teal!4,colframe=teal!35!black,title={Before you move on}]
What does \te{వలసపోవు} mean? (\textquotedblleft{}To migrate\textquotedblright{}.) Say it once more.
\end{tcolorbox}
